\section{热力学基础}
\sourceproblem{选择1}{58}{64}
\problemtext 固定量理想气体由等压的两个平衡态A变至B，图示末体积较大，问必然变化。
\workings 由 $pV=\nu RT$ 得 $T_B>T_A$，固定组成理想气体内能随温度增加，选B。
\crosscheck 仅端点压强相等，不代表途中等压，不能据此断定功或热量。
\sourceproblem{选择2}{58}{64}
\problemtext 单原子理想气体等压膨胀，求对外功与吸热之比。
\workings $W=\nu R\Delta T,\ Q=\nu C_p\Delta T$，$C_p=5R/2$，故 $W/Q=2/5$，选C。
\crosscheck 内能增加占吸热的3/5；取经典单原子定热容模型。
\sourceproblem{选择3}{58}{64}
\problemtext 常温氧气、氦气初始压强体积温度相等，分别等压吸收相同热量，比较功。
\workings 刚性氧气 $C_p=7R/2$，氦气 $C_p=5R/2$。等压 $W=RQ/C_p$，故 $W_{\rm O_2}:W_{\rm He}=5:7$，选C。
\crosscheck 气体种类经原页放大确认；若振动自由度显著激发，定热容近似需修正。
\sourceproblem{选择4}{58}{64}
\problemtext 比较等容加热、等温压缩、等压压缩及绝热压缩的热力学陈述。
\workings 取系统对外功为正。绝热压缩 $Q=0,W<0$，故 $\Delta U=-W>0$，温度升高且体积减小，压强升高，选D。等容加热内能增加；等温压缩放热；等压压缩内能减少。
\crosscheck 采用固定量理想气体及题设简单压缩过程。
\sourceproblem{选择5}{58}{64}
\problemtext 同一初态的2 mol理想氧气分别等温、等压、等容吸收400 J，比较功和内能增量。
\workings 等温 $W=400$ J、$\Delta U=0$；等压 $W=2Q/7$、$\Delta U=5Q/7$；等容 $W=0,\Delta U=400$ J。因此等温功最大，等容内能增加最多，选D。
\crosscheck 等压热容按刚性双原子模型；即使热容变化，定性排序仍成立。
\sourceproblem{选择6}{58}{64}
\problemtext 图2理想气体先等容增压再等压膨胀，判断热量及内能变化。
\workings 两段均使 $pV$ 增大，温度升高，$\Delta U>0$；第一段无体积功，第二段对外功为正，故 $Q=\Delta U+W>0$，选A。
\crosscheck 图中过程方向为a至b至c；端点温度由乘积而非比值确定。
\sourceproblem{选择7}{59}{65}
\problemtext 图3从共同初态A沿三路径膨胀至同温的B、C、D，AC绝热，判断AB和AD吸放热。
\workings 三末态内能相同，AC有 $W_C=-\Delta U$。图中AB路径位于AC下方且终体积较小，AD位于AC上方且终体积较大，故 $W_B<W_C<W_D$，所以 $Q_B<0<Q_D$，选B。
\crosscheck 功为曲线下有向面积；必须比较完整路径而非仅末态压强。
\sourceproblem{选择8}{59}{65}
\problemtext 从同一初态压缩到初体积一半，分别等压、等温、绝热，比较外界对气体做功。
\workings 可逆绝热 $p=p_0(V_0/V)^\gamma$，等温 $p=p_0V_0/V$，等压 $p=p_0$。当 $V<V_0,\gamma>1$，压强依次为绝热、等温、等压从大到小，积分外界功得到相同排序，选A。
\crosscheck 此处比较的是图示准静态路径的外界功，均为正值。
\sourceproblem{选择9}{59}{65}
\problemtext 根据图5的V--T循环判定内能增加的过程。
\workings 温度是横轴；A至B温度增加，B至C温度降低，C至A温度恒定，因此仅A至B内能增加，选A。
\crosscheck 固定组成理想气体内能仅随温度变化；纵轴体积变化不直接决定内能。
\sourceproblem{选择10}{59}{65}
\problemtext 判断热力学第二定律相关陈述。
\workings 理想气体向真空自由绝热膨胀，内能和温度不变，而 $\Delta S=\nu R\ln(V_f/V_i)>0$，选C。
\crosscheck 单次等温膨胀可将吸热全部转为功，故禁止单热源全热转功须附加循环且无其他影响条件；制冷机可借外功使热由低温流向高温。
\sourceproblem{填空1}{59}{65}
\problemtext 容器内分子数 $N$、体积 $V$、质量 $m$、摩尔质量 $M$，求数密度、物质的量及与阿伏伽德罗常数关系。
\workings $n=N/V,\ \nu=m/M,\ N=\nu N_A=mN_A/M$。
\crosscheck 数密度单位为 $\mathrm{m^{-3}}$，物质的量单位为mol。
\sourceproblem{填空2}{59}{65}
\problemtext 系统吸热826 J并对外做功500 J，求内能改变。
\workings $\Delta U=Q-W=326$ J。
\crosscheck 采用吸热和系统对外功为正。
\sourceproblem{填空3}{59}{65}
\problemtext 理想气体等压对外做功200 J，分别求单原子和刚性双原子气体吸热。
\workings $\nu R\Delta T=200$ J。$Q=\nu C_p\Delta T$，故单原子为500 J，双原子为700 J。
\crosscheck 相应内能增量为300和500 J，均满足第一定律。
\sourceproblem{填空4}{59}{65}
\problemtext 10 mol单原子理想气体受外界压缩功209 J，温度升高1 K，求内能增加及外界传热（续PDF66）。
\workings $\Delta U=3\nu R\Delta T/2=124.65$ J（$R=8.31$），系统对外功 $W=-209$ J，故 $Q=\Delta U+W=-84.35$ J，即放热84.35 J。
\crosscheck 压缩做功不意味着气体一定吸热；能量输入大于内能增量。
\sourceproblem{填空5}{60}{66}
\problemtext 图6矩形顺时针循环，压强10至40 Pa、体积1至4立方米，求净吸热。
\workings 循环 $\Delta U=0$，净功为面积 $(40-10)(4-1)=90$ J，故净吸热90 J。
\crosscheck 顺时针循环净功为正；等容段无体积功。
\sourceproblem{填空6}{60}{66}
\problemtext 图7从a经acb到b吸热80 J、做功30 J；经adb做功10 J；由b沿ba返回时外界做功20 J，求后两段热量。
\workings $U_b-U_a=80-30=50$ J。路径adb吸热 $50+10=60$ J。返回时 $Q=-50-20=-70$ J，即放热70 J。
\crosscheck 走acb后返回的循环净功10 J、净吸热10 J。
\sourceproblem{填空7}{60}{66}
\problemtext 图8为顺时针椭圆状循环，压强轴从1至3 Pa、体积轴从1至3立方米，求一周净功。
\workings 若图示为标准椭圆（以坐标单位计半轴均为1），净功等于椭圆面积 $\pi(1\,\mathrm{Pa})(1\,\mathrm{m^3})=\pi$ J。
\crosscheck 图形未另给解析方程；数值依标准椭圆读图假设，若只是示意曲线则只能确定净功为正及面积积分形式。
\sourceproblem{填空8}{60}{66}
\problemtext 图9沿直线从坐标(2,4)至(4,2)，求气体对外功与内能改变。
\workings 令每格体积为 $s_V$、每格压强为 $s_p$，则 $W=(4+2)(4-2)s_ps_V/2=6s_ps_V$。两端 $pV=8s_ps_V$，故 $\Delta U=0$。若按常见图示单位取 $s_p=10^3$ Pa、$s_V=10^{-3}$ 立方米，得 $W=6$ J。
\crosscheck 原图两个轴均标注乘 $10^3$，若严格将其读成坐标数值等于物理量乘 $10^3$，则两尺度均为 $10^{-3}$，功为 $6\times10^{-6}$ J；保留单位歧义标记，不把条件数值认定为唯一答案。直线路径面积及零内能差不受此歧义影响。
\sourceproblem{填空9}{60}{66}
\problemtext 两卡诺热机具有同一低温热源、不同高温热源，p--V循环面积相等，比较净功和热效率。
\workings 两个正向热机循环的净功均等于面积，故相等；$\eta=1-T_L/T_H$，固定正的 $T_L$ 时随 $T_H$ 严格增加，因此效率不同。
\crosscheck 等净功不意味着吸热量相同；热源温度须用开尔文。
